% ============================================================================
%  ESTADÍSTICA DESCRIPTIVA — Sesión 14
%  Semana 14: 9 al 14 de noviembre de 2026
%  3.3 Clasificación de variables y construcción de tablas en un caso completo
%  3.4 Cálculo de medidas y elaboración de gráficos con herramientas
%      tecnológicas  (R base)
%  Duración prevista: 3 horas
% ============================================================================
\documentclass[aspectratio=169,10pt]{beamer}

\usepackage[utf8]{inputenc}
\IfFileExists{spanish.ldf}{\usepackage[spanish,es-noshorthands]{babel}}{}
\usepackage{amsmath,amssymb}
\usepackage{booktabs}
\usepackage{tikz}
\usepackage{listings}

\usetheme{Uniminuto}

\uniasignatura{Estadística descriptiva}
\uniprofesor{Andrés Felipe Puerta Velez}
\unicorreo{andres.puerta-v@uniminuto.edu.co}
\uninrc{CBASBL031 - CBASBL141}
\unisemana{Sesión 14 — Semana 14 · 9 al 14 de noviembre de 2026}

\newcommand{\listatiny}{\setbeamerfont{itemize/enumerate body}{size=\tiny}\tiny}
\newcommand{\fx}[1]{{\ttfamily\color{uniAzul}#1}}

% ------------------------------------------------------------- CÓDIGO R ----
\lstdefinestyle{Rcode}{
  language=R,
  basicstyle=\ttfamily\scriptsize,
  keywordstyle=\color{uniAzul}\bfseries,
  commentstyle=\color{uniGris}\itshape,
  stringstyle=\color{uniTitulo},
  showstringspaces=false,
  numbers=none,
  breaklines=true,
  columns=fullflexible,
  backgroundcolor=\color{uniCremaClaro},
  frame=leftline, framerule=1.2pt, rulecolor=\color{uniNaranja},
  xleftmargin=0.35cm, framexleftmargin=0.35cm,
  framexrightmargin=0.15cm, framextopmargin=2pt, framexbottommargin=2pt,
  aboveskip=0.15cm, belowskip=0.15cm
}
\lstdefinestyle{Rmini}{
  language=R,
  basicstyle=\ttfamily\tiny,
  keywordstyle=\color{uniAzul}\bfseries,
  commentstyle=\color{uniGris}\itshape,
  stringstyle=\color{uniTitulo},
  showstringspaces=false, numbers=none, breaklines=true,
  columns=fullflexible,
  backgroundcolor=\color{uniCremaClaro},
  frame=leftline, framerule=1.2pt, rulecolor=\color{uniNaranja},
  xleftmargin=0.35cm, framexleftmargin=0.35cm,
  framexrightmargin=0.15cm, framextopmargin=2pt, framexbottommargin=2pt,
  aboveskip=0.12cm, belowskip=0.12cm
}
\lstdefinestyle{Rout}{
  basicstyle=\ttfamily\scriptsize\color{black},
  showstringspaces=false, numbers=none, breaklines=true,
  columns=fullflexible,
  backgroundcolor=\color{uniGris!10},
  frame=leftline, framerule=1.2pt, rulecolor=\color{uniGris},
  xleftmargin=0.35cm, framexleftmargin=0.35cm,
  framexrightmargin=0.15cm, framextopmargin=2pt, framexbottommargin=2pt,
  aboveskip=0.15cm, belowskip=0.15cm
}

\begin{document}

\portada

\begin{frame}{Orden del día}
  \begin{enumerate}
    \item Llamado a lista y revisión de los ejercicios propuestos
    \item \azul{Tema 3.3: un caso completo, de los datos crudos a la tabla}
    \item Primeros pasos en \azul{R}: objetos, vectores y \fx{data.frame}
    \item \azul{Tema 3.4: medidas, tablas y gráficos con R base}
    \item Comparación de grupos y errores frecuentes
    \item Taller en computador y cierre
  \end{enumerate}

  \vspace{0.25cm}
  {\footnotesize\color{uniGris} Sesión de 3 horas. Traiga computador si puede;
  si no, se trabaja en parejas. Usamos \destacar{R base}: nada de lo que
  veremos hoy exige instalar paquetes adicionales.}
\end{frame}

\begin{frame}{Por qué R y no una hoja de cálculo}
  \scriptsize
  \begin{columns}[T,onlytextwidth]
    \begin{column}{0.48\textwidth}
      \begin{block}{Lo que gana}
        \tiny
        \begin{itemize}\setlength{\itemsep}{-0.05cm}
          \item El análisis queda \azul{escrito}: se puede repetir, revisar y
                corregir.
          \item Si cambian los datos, se vuelve a correr el mismo guion.
          \item Otra persona puede \azul{reproducir} exactamente su
                resultado.
          \item Es gratuito y es el estándar en analítica.
        \end{itemize}
      \end{block}
    \end{column}
    \begin{column}{0.48\textwidth}
      \begin{alertblock}{Lo que cuesta}
        \tiny
        \begin{itemize}\setlength{\itemsep}{-0.05cm}
          \item Hay que escribir, no hacer clic.
          \item Los errores de sintaxis detienen todo (mayúsculas, comas,
                paréntesis).
          \item La curva de entrada es de unas pocas horas.
        \end{itemize}
      \end{alertblock}
    \end{column}
  \end{columns}

  \vspace{0.1cm}
  \begin{exampleblock}{La razón de fondo}
    \tiny
    En una hoja de cálculo, el análisis \destacar{está en los clics que usted
    hizo} y nadie los ve. En R, el análisis \destacar{es el archivo}: cualquiera
    puede leerlo, auditarlo y volverlo a correr dentro de un año.
  \end{exampleblock}
\end{frame}

\begin{frame}[fragile]{Lo mínimo para empezar}
  \footnotesize
  \begin{columns}[T,onlytextwidth]
    \begin{column}{0.50\textwidth}
      \begin{block}{Cuatro ideas}
        \tiny
        \begin{itemize}\setlength{\itemsep}{-0.05cm}
          \item \fx{<-} guarda un valor en un objeto.
          \item \fx{c()} arma un vector (una columna de datos).
          \item Una \azul{función} se llama con paréntesis:
                \fx{mean(x)}.
          \item \fx{\#} inicia un comentario: R lo ignora.
        \end{itemize}
      \end{block}
    \end{column}
    \begin{column}{0.46\textwidth}
      \begin{alertblock}{Dónde se escribe}
        \tiny
        En RStudio: el \azul{Script} (arriba a la izquierda) se guarda; la
        \azul{Consola} (abajo) solo ejecuta. \destacar{Escriba siempre en el
        script} y ejecute con Ctrl+Enter.
      \end{alertblock}
    \end{column}
  \end{columns}

  \vspace{0.05cm}
\begin{lstlisting}[style=Rcode]
# Un objeto con tres numeros
x <- c(4, 7, 10)
mean(x)       # promedio
length(x)     # cuantos datos hay
\end{lstlisting}
\begin{lstlisting}[style=Rout]
[1] 7
[1] 3
\end{lstlisting}
\end{frame}

\seccion{3.3 Caso completo: Ferretería El Tornillo}

\begin{frame}{Los datos crudos}
  \footnotesize
  Valor de compra, en miles de pesos, de 50 clientes de octubre (ya ordenados):

  \vspace{0.15cm}
  \begin{beamercolorbox}[wd=\textwidth,sep=0.12cm]{block body}
    \centering\scriptsize\ttfamily
    10 \ 13 \ 16 \ 19 \ 21 \ 22 \ 23 \ 24 \ 25 \ 26 \ 27 \ 28 \ 28 \ 29 \ 29 \\
    30 \ 31 \ 32 \ 33 \ 34 \ 35 \ 35 \ 36 \ 36 \ 37 \ 38 \ 38 \ 39 \ 39 \ 41 \\
    42 \ 43 \ 44 \ 45 \ 46 \ 47 \ 48 \ 49 \ 51 \ 53 \ 54 \ 56 \ 57 \ 58 \ 61 \\
    63 \ 66 \ 68 \ 72 \ 80
  \end{beamercolorbox}

  \vspace{0.15cm}
  \begin{alertblock}{Antes de calcular nada}
    \scriptsize
    \destacar{Contar.} $n = 50$. \quad \destacar{Mirar los extremos.}
    $x_{\min}=10$, $x_{\max}=80$. \quad \destacar{Buscar rarezas.} ¿Hay ceros,
    negativos, datos faltantes o valores imposibles? Aquí no, pero ese chequeo
    es obligatorio en datos reales.
  \end{alertblock}
\end{frame}

\begin{frame}{Paso 1 — clasificar las variables}
  \scriptsize
  \begin{center}
  \begin{tabular}{@{}p{2.6cm}p{3.0cm}p{2.5cm}p{3.4cm}@{}}
    \toprule
    \textbf{Variable} & \textbf{Registro} & \textbf{Tipo y escala} &
    \textbf{Cómo se guarda en R} \\
    \midrule
    Valor de compra & Miles de pesos & Continua, razón &
      \fx{numeric} \\
    \addlinespace
    Zona del cliente & Norte / Centro / Sur & Nominal &
      \fx{factor} \\
    \addlinespace
    Forma de pago & Efectivo / Tarjeta / Crédito & Nominal &
      \fx{factor} \\
    \addlinespace
    Nivel de fidelidad & Nuevo / Ocasional / Frecuente & Ordinal &
      \fx{factor(ordered = TRUE)} \\
    \bottomrule
  \end{tabular}
  \end{center}

  \vspace{0.1cm}
  \begin{alertblock}{Por qué importa el tipo en R}
    \scriptsize
    Si una variable cualitativa queda como \fx{character}, R la ordena
    \azul{alfabéticamente}. Si se declara \fx{factor} con niveles en el orden
    correcto, las tablas y los gráficos respetan el orden de la escala. Es la
    traducción exacta de lo que vimos en la sesión 2.
  \end{alertblock}
\end{frame}

\begin{frame}[fragile]{Paso 2 — cargar los datos en R}
  \footnotesize
\begin{lstlisting}[style=Rcode]
# Opcion A: escribir el vector directamente
valor <- c(10,13,16,19,21,22,23,24,25,26,27,28,28,29,29,
           30,31,32,33,34,35,35,36,36,37,38,38,39,39,41,
           42,43,44,45,46,47,48,49,51,53,54,56,57,58,61,
           63,66,68,72,80)

# Opcion B: leer un archivo guardado como CSV
# datos <- read.csv("tornillo.csv")
# valor <- datos$valor

length(valor)   # verificar que son 50
summary(valor)  # primer vistazo
\end{lstlisting}
\begin{lstlisting}[style=Rout]
[1] 50
   Min. 1st Qu.  Median    Mean 3rd Qu.    Max.
  10.00   28.25   37.50   39.54   48.75   80.00
\end{lstlisting}

  \vspace{-0.1cm}
  {\tiny\color{uniGris} \fx{summary()} es el primer comando que se corre
  siempre: en una línea entrega el resumen de cinco números y la media.}
\end{frame}

\seccion{3.4 Medidas, tablas y gráficos con R base}

\begin{frame}{Las funciones que reemplazan todo el cálculo a mano}
  \tiny
  \begin{center}
  \begin{tabular}{@{}p{3.2cm}p{5.0cm}p{3.4cm}@{}}
    \toprule
    \textbf{Medida} & \textbf{Función en R} & \textbf{Resultado} \\
    \midrule
    Cantidad de datos & \fx{length(valor)} & 50 \\
    Media & \fx{mean(valor)} & 39{,}54 \\
    Mediana & \fx{median(valor)} & 37{,}50 \\
    Mínimo y máximo & \fx{min(valor)} \ \fx{max(valor)} & 10 \ y \ 80 \\
    Rango & \fx{diff(range(valor))} & 70 \\
    Cuartiles & \fx{quantile(valor)} & 28{,}25 \ y \ 48{,}75 \\
    Percentil 90 & \fx{quantile(valor, 0.90)} & 61{,}20 \\
    Rango intercuartílico & \fx{IQR(valor)} & 20{,}50 \\
    Varianza (muestra) & \fx{var(valor)} & 247{,}15 \\
    Desviación (muestra) & \fx{sd(valor)} & 15{,}72 \\
    Coeficiente de variación & \fx{sd(valor)/mean(valor)*100} & 39{,}76\% \\
    Resumen completo & \fx{summary(valor)} & los cinco números y la media \\
    \bottomrule
  \end{tabular}
  \end{center}

  \vspace{0.05cm}
  {\tiny\color{uniGris} R \azul{no} trae una función de moda ni de asimetría en
  la instalación base: las construimos en las dos láminas siguientes.
  \fx{var()} y \fx{sd()} dividen entre $n-1$, es decir, son las de
  \destacar{muestra}.}
\end{frame}

\begin{frame}[fragile]{Moda, asimetría y curtosis a mano en R}
  \footnotesize
\begin{lstlisting}[style=Rcode]
# Moda: el valor mas repetido
tab <- table(valor)
names(tab)[tab == max(tab)]

# Asimetria de Pearson (mediana)
As2 <- 3 * (mean(valor) - median(valor)) / sd(valor)

# Fisher: g1 y g2 (exceso de curtosis), con sigma poblacional
n  <- length(valor)
mu <- mean(valor)
sg <- sqrt(sum((valor - mu)^2) / n)
g1 <- sum((valor - mu)^3) / n / sg^3
g2 <- sum((valor - mu)^4) / n / sg^4 - 3
round(c(As2 = As2, g1 = g1, g2 = g2), 3)
\end{lstlisting}
\begin{lstlisting}[style=Rout]
[1] "28" "29" "35" "36" "38" "39"
   As2     g1     g2
 0.389  0.465 -0.242
\end{lstlisting}

  \vspace{-0.1cm}
  {\tiny\color{uniGris} Seis valores empatados con dos repeticiones: en datos
  continuos la moda \destacar{casi nunca es informativa}, y por eso se agrupa.}
\end{frame}

\begin{frame}[fragile]{La tabla de frecuencias con \fx{cut()} y \fx{table()}}
  \footnotesize
\begin{lstlisting}[style=Rcode]
limites <- seq(10, 80, by = 10)          # 10,20,...,80
clases  <- cut(valor, breaks = limites,
               right = FALSE,            # intervalos [a, b)
               include.lowest = TRUE)    # el 80 entra en la ultima
fi <- table(clases)
fi
\end{lstlisting}
\begin{lstlisting}[style=Rout]
clases
[10,20) [20,30) [30,40) [40,50) [50,60) [60,70) [70,80]
      4      11      14       9       6       4       2
\end{lstlisting}

  \vspace{0.05cm}
  \begin{alertblock}{Los dos argumentos que hay que entender}
    \tiny
    \fx{right = FALSE} produce intervalos \destacar{cerrados por izquierda y
    abiertos por derecha}, que es la convención del curso. \quad
    \fx{include.lowest = TRUE} cierra el \destacar{último} intervalo para que
    el máximo no quede fuera. Sin estos dos argumentos, R usa $(a, b]$ y la
    tabla no coincide con la que hicimos a mano.
  \end{alertblock}
\end{frame}

\begin{frame}[fragile]{La tabla completa, con todas las frecuencias}
  \footnotesize
\begin{lstlisting}[style=Rcode]
Fi <- cumsum(fi)                  # absoluta acumulada
hi <- prop.table(fi)              # relativa
Hi <- cumsum(hi)                  # relativa acumulada
xi <- head(limites, -1) + 5       # marca de clase

tabla <- data.frame(xi, fi = as.vector(fi), Fi = as.vector(Fi),
                    hi = round(as.vector(hi), 3),
                    Hi = round(as.vector(Hi), 3))
tabla
\end{lstlisting}
\begin{lstlisting}[style=Rout]
  xi fi Fi    hi    Hi
1 15  4  4 0.080 0.080
2 25 11 15 0.220 0.300
3 35 14 29 0.280 0.580
4 45  9 38 0.180 0.760
5 55  6 44 0.120 0.880
6 65  4 48 0.080 0.960
7 75  2 50 0.040 1.000
\end{lstlisting}

  \vspace{-0.12cm}
  {\tiny\color{uniGris} Es \destacar{exactamente} la tabla de la sesión 11,
  pero en seis líneas y sin una sola cuenta a mano.}
\end{frame}

\begin{frame}[fragile]{Variables cualitativas: \fx{factor} y \fx{table}}
  \footnotesize
\begin{lstlisting}[style=Rcode]
zona <- factor(zona)                       # nominal
fidelidad <- factor(fidelidad,
                    levels = c("Nuevo", "Ocasional", "Frecuente"),
                    ordered = TRUE)        # ordinal: el orden manda

table(zona)                                # frecuencias absolutas
round(prop.table(table(zona)) * 100, 1)    # porcentajes
cumsum(table(fidelidad))                   # acumulada: solo si es ordinal
\end{lstlisting}

  \vspace{0.05cm}
  \begin{alertblock}{El detalle que decide todo}
    \scriptsize
    Al declarar \fx{levels = c("Nuevo","Ocasional","Frecuente")} le estamos
    diciendo a R cuál es el \destacar{orden natural} de la escala. Sin eso,
    R ordenaría alfabéticamente —Frecuente, Nuevo, Ocasional— y tanto la
    acumulada como el gráfico de barras quedarían sin sentido.
  \end{alertblock}

  \vspace{0.05cm}
  {\tiny\color{uniGris} Para cruzar dos variables cualitativas:
  \fx{table(zona, pago)} devuelve la tabla de doble entrada, y
  \fx{addmargins()} le agrega los totales.}
\end{frame}

\begin{frame}[fragile]{Los gráficos, uno por línea}
  \footnotesize
\begin{lstlisting}[style=Rmini]
# Cualitativa: barras separadas y circular
barplot(table(zona), col = "steelblue",
        main = "Clientes por zona", ylab = "Frecuencia")
pie(table(zona), main = "Participacion por zona")

# Cuantitativa: histograma con NUESTROS limites de clase
hist(valor, breaks = limites, right = FALSE,
     main = "Valor de compra de 50 clientes",
     xlab = "Miles de pesos", ylab = "Frecuencia", col = "steelblue")

# Poligono de frecuencias sobre el histograma
lines(c(5, xi, 85), c(0, as.vector(fi), 0), type = "b", col = "tomato")

# Ojiva: acumulada sobre los limites superiores
plot(limites, c(0, Fi), type = "b",
     main = "Ojiva", xlab = "Miles de pesos", ylab = "Fi acumulada")

# Diagrama de caja
boxplot(valor, horizontal = TRUE, col = "steelblue",
        main = "Valor de compra", xlab = "Miles de pesos")
\end{lstlisting}

  \vspace{-0.12cm}
  {\tiny\color{uniGris} \fx{hist()} con \fx{breaks = limites} obliga a R a usar
  las clases del curso; sin ese argumento, R elige las suyas con la regla de
  Sturges y el gráfico no coincide con la tabla.}
\end{frame}

\begin{frame}[fragile]{Comparar grupos en dos líneas}
  \footnotesize
\begin{lstlisting}[style=Rcode]
# Una medida por grupo
tapply(valor, zona, mean)
tapply(valor, zona, median)
round(tapply(valor, zona, sd) / tapply(valor, zona, mean) * 100, 1)  # CV

# Cajas paralelas: la formula se lee "valor SEGUN zona"
boxplot(valor ~ zona, col = "steelblue",
        main = "Valor de compra por zona",
        xlab = "Zona", ylab = "Miles de pesos")

# Los cinco numeros y los atipicos que usa el grafico
boxplot.stats(valor)$stats   # bigote inf, Q1, mediana, Q3, bigote sup
boxplot.stats(valor)$out     # valores atipicos
\end{lstlisting}
\begin{lstlisting}[style=Rout]
[1] 10 28 37.5 49 80
numeric(0)
\end{lstlisting}

  \vspace{-0.12cm}
  {\tiny\color{uniGris} \fx{numeric(0)} significa ``ningún valor atípico'',
  igual que concluimos a mano en la sesión 11. Esta lámina es toda la
  comparación de grupos de la sesión 13, hecha en R.}
\end{frame}

\begin{frame}{¿Por qué R no da exactamente lo mismo que a mano?}
  \tiny
  \begin{center}
  \begin{tabular}{@{}lrrl@{}}
    \toprule
    \textbf{Medida} & \textbf{A mano (agrupado)} & \textbf{R (datos crudos)} &
    \textbf{Por qué difieren} \\
    \midrule
    Media    & 39{,}40 & 39{,}54 & La marca de clase reemplaza al dato real \\
    Mediana  & 37{,}14 & 37{,}50 & La fórmula agrupada interpola dentro de la clase \\
    Moda     & 33{,}75 & varias & Czuber estima; \fx{table()} da los valores repetidos \\
    $Q_1$    & 27{,}73 & 28{,}25 & Distinta convención de posición \\
    $Q_3$    & 49{,}44 & 48{,}75 & Distinta convención de posición \\
    $s$      & 15{,}54 & 15{,}72 & Agrupamiento \\
    \bottomrule
  \end{tabular}
  \end{center}

  \vspace{0.1cm}
  \begin{alertblock}{Cuál está bien}
    \scriptsize
    \destacar{Los dos.} El cálculo agrupado es una \azul{aproximación} que se
    usa cuando solo se dispone de la tabla; R usa los datos originales y es más
    exacto. Lo que \destacar{nunca} se hace es mezclarlos en un mismo informe
    sin decir cuál se usó.
  \end{alertblock}

  \vspace{0.05cm}
  {\tiny\color{uniGris} \fx{quantile()} tiene \azul{nueve} métodos distintos
  (\fx{type = 1..9}); el predeterminado es el 7. \fx{fivenum()} y
  \fx{boxplot()} usan otra convención (bisagras de Tukey), y por eso
  \fx{quantile(valor)} da 28,25 mientras \fx{boxplot.stats()} da 28.
  Diferencias así son normales y no son un error.}
\end{frame}

\begin{frame}{Errores frecuentes en R}
  \scriptsize
  \begin{center}
  \begin{tabular}{@{}p{5.0cm}p{7.0cm}@{}}
    \toprule
    \textbf{Error} & \textbf{Qué pasa y cómo se corrige} \\
    \midrule
    \fx{Error: object 'valor' not found} &
      No se ejecutó la línea que lo crea, o está mal escrito. R distingue
      mayúsculas: \fx{Valor} $\ne$ \fx{valor} \\
    \addlinespace
    La consola muestra \fx{+} y no responde &
      Falta cerrar un paréntesis o una comilla. Pulse \destacar{Esc} \\
    \addlinespace
    \fx{mean()} devuelve \fx{NA} &
      Hay datos faltantes: use \fx{mean(valor, na.rm = TRUE)} \\
    \addlinespace
    La tabla no coincide con la de clase &
      Falta \fx{right = FALSE} o \fx{include.lowest = TRUE} en \fx{cut()} \\
    \addlinespace
    El gráfico de barras sale desordenado &
      La variable es \fx{character}: conviértala en \fx{factor} con
      \fx{levels} \\
    \addlinespace
    \fx{sd()} distinto al esperado &
      \fx{sd()} es de \azul{muestra} ($n-1$); la poblacional se calcula a mano \\
    \bottomrule
  \end{tabular}
  \end{center}

  \vspace{0.05cm}
  {\tiny\color{uniGris} Control final infalible: \fx{sum(fi)} debe dar
  \destacar{exactamente} \fx{length(valor)}. Si no coincide, hay datos fuera de
  los límites de \fx{cut()}.}
\end{frame}

\begin{frame}{Taller en computador — parte 1}
  \begin{alertblock}{En parejas (30 minutos)}
    \scriptsize
    Abra RStudio, cree un script nuevo y guárdelo como
    \fx{tornillo.R}. Todo debe quedar \destacar{escrito en el script}, no en la
    consola.

    \vspace{0.2cm}
    \begin{enumerate}\setlength{\itemsep}{-0.07cm}
      \item Cargue los 50 datos con \fx{c()} y verifique \fx{length()}.
      \item Corra \fx{summary()} e interprete su salida en una frase.
      \item Calcule media, mediana, cuartiles, \fx{IQR}, \fx{var}, \fx{sd} y el
            $CV$.
      \item Construya \fx{clases} con \fx{cut()} y obtenga \fx{fi}.
      \item Agregue \fx{Fi}, \fx{hi}, \fx{Hi} y \fx{xi}, y arme el
            \fx{data.frame}.
      \item Verifique que \fx{sum(fi) == 50} y que el último \fx{Hi} es 1.
      \item Compare sus resultados con los calculados a mano en la sesión 11 y
            explique cada diferencia.
    \end{enumerate}
  \end{alertblock}
\end{frame}

\begin{frame}{Taller en computador — parte 2}
  \begin{alertblock}{En parejas (25 minutos)}
    \scriptsize
    \begin{enumerate}\setlength{\itemsep}{-0.07cm}
      \item Haga el \azul{histograma} con \fx{breaks = limites} y
            \fx{right = FALSE}.
      \item Superponga el \azul{polígono} con \fx{lines()}, cerrado sobre el
            eje.
      \item Dibuje la \azul{ojiva} con \fx{plot()}.
      \item Dibuje el \azul{diagrama de caja} y compárelo con el que hizo a
            mano en la sesión 11.
      \item Cree una variable \fx{zona} con 50 valores inventados
            (Norte/Centro/Sur), conviértala en \fx{factor} y haga
            \fx{table(zona)} y su gráfico de barras.
      \item Calcule \fx{tapply(valor, zona, mean)} y haga
            \fx{boxplot(valor \~{} zona)}. Escriba \destacar{dos frases} de
            interpretación.
    \end{enumerate}

    \vspace{0.1cm}
    Al cierre, dos parejas proyectan su script y el curso revisa títulos, ejes
    y coherencia de los gráficos.
  \end{alertblock}
\end{frame}

\begin{frame}[fragile]{El script completo, para guardar}
  \footnotesize
\begin{lstlisting}[style=Rmini]
valor   <- c(10,13,16,19,21,22,23,24,25,26,27,28,28,29,29,30,31,32,33,34,
             35,35,36,36,37,38,38,39,39,41,42,43,44,45,46,47,48,49,51,53,
             54,56,57,58,61,63,66,68,72,80)
limites <- seq(10, 80, by = 10)
clases  <- cut(valor, breaks = limites, right = FALSE, include.lowest = TRUE)

fi <- table(clases); Fi <- cumsum(fi)
hi <- prop.table(fi); Hi <- cumsum(hi)
xi <- head(limites, -1) + 5
tabla <- data.frame(xi, fi = as.vector(fi), Fi = as.vector(Fi),
                    hi = round(as.vector(hi), 3), Hi = round(as.vector(Hi), 3))

medidas <- c(n = length(valor), media = mean(valor), mediana = median(valor),
             Q1 = quantile(valor, .25), Q3 = quantile(valor, .75),
             IQR = IQR(valor), s = sd(valor),
             CV = sd(valor) / mean(valor) * 100)

hist(valor, breaks = limites, right = FALSE, col = "steelblue",
     main = "Valor de compra", xlab = "Miles de pesos")
boxplot(valor, horizontal = TRUE, col = "steelblue")
\end{lstlisting}

  \vspace{-0.12cm}
  {\tiny\color{uniGris} Diecisiete líneas reemplazan tres sesiones de cálculo a
  mano. \destacar{Pero solo tienen sentido si usted sabe qué está calculando}:
  esa es la diferencia entre las sesiones 3 a 11 y esta.}
\end{frame}

\begin{frame}{Lista de verificación antes de entregar}
  \scriptsize
  \begin{columns}[T,onlytextwidth]
    \begin{column}{0.48\textwidth}
      \begin{block}{Script y cálculos}
        \tiny\setbeamerfont{itemize/enumerate body}{size=\tiny}
        \begin{itemize}\setlength{\itemsep}{-0.05cm}
          \item \fx{length(valor)} coincide con el $n$ declarado
          \item \fx{sum(fi)} $=$ \fx{length(valor)} y el último \fx{Hi} es 1
          \item Se usó \fx{sd()} sabiendo que es de muestra
          \item Ninguna cifra está digitada a mano en el informe
          \item El script corre completo \azul{de arriba abajo} sin errores
        \end{itemize}
      \end{block}
    \end{column}
    \begin{column}{0.48\textwidth}
      \begin{block}{Tablas y gráficos}
        \tiny\setbeamerfont{itemize/enumerate body}{size=\tiny}
        \begin{itemize}\setlength{\itemsep}{-0.05cm}
          \item Cada tabla y figura está numerada y titulada
          \item \fx{main}, \fx{xlab} y \fx{ylab} escritos con unidades
          \item El histograma usa \fx{breaks = limites}
          \item El polígono está cerrado sobre el eje
          \item Cada figura se menciona en el texto
        \end{itemize}
      \end{block}
    \end{column}
  \end{columns}

  \vspace{0.05cm}
  \begin{alertblock}{La prueba definitiva}
    \tiny
    Cierre R, ábralo de nuevo y corra el script \destacar{desde la primera
    línea}. Si todo se reproduce sin tocar nada, el trabajo está bien hecho.
  \end{alertblock}
\end{frame}

\begin{frame}{Cierre de la sesión}
  \begin{block}{Lo que debe quedar claro hoy}
    \scriptsize
    \begin{itemize}\setlength{\itemsep}{0.02cm}
      \item En R el análisis \azul{es un archivo}: se guarda, se audita y se
            repite.
      \item \fx{summary()}, \fx{mean()}, \fx{median()}, \fx{quantile()},
            \fx{IQR()}, \fx{var()} y \fx{sd()} cubren toda la Unidad 2.
      \item \fx{cut()} + \fx{table()} + \fx{cumsum()} + \fx{prop.table()}
            construyen la tabla de la Unidad 1.
      \item \fx{hist()}, \fx{barplot()}, \fx{plot()} y \fx{boxplot()} dan todos
            los gráficos del curso.
      \item Los resultados con datos crudos \destacar{difieren un poco} de los
            agrupados, y hay que declarar cuál se usó.
      \item R no piensa: el criterio sobre qué medida y qué gráfico usar sigue
            siendo suyo.
    \end{itemize}
  \end{block}

  \vspace{0.08cm}
  \begin{alertblock}{Trabajo independiente}
    \tiny
    Terminar el script del taller. En la sesión 15 trabajamos la
    \destacar{coherencia entre análisis y conclusiones} y la sustentación oral,
    y hacemos el repaso para el Examen 3.
  \end{alertblock}
\end{frame}

\begin{frame}{Ejercicios propuestos}
  \vspace{-0.1cm}
  \begin{beamercolorbox}[wd=\textwidth,sep=0.10cm]{block body}
    \centering\scriptsize
    \textbf{1--8:} funciones. \quad \textbf{9--14:} tablas y gráficos. \quad
    \textbf{15--20:} criterio y verificación.
  \end{beamercolorbox}
  \listatiny
  \begin{columns}[T,onlytextwidth]
    \begin{column}{0.47\textwidth}
      \begin{enumerate}\setlength{\itemsep}{0.16cm}
        \item Función para la media de \fx{valor}
        \item Función para la mediana
        \item Función para el tercer cuartil
        \item Función para el percentil 40
        \item ¿\fx{sd()} es de muestra o de población?
        \item ¿Cómo se calcula el $CV$ en R?
        \item Función para contar los clientes de la zona Norte
        \item ¿Qué hace \fx{summary(valor)}?
        \item ¿Qué argumentos de \fx{cut()} dan intervalos $[a,b)$?
        \item ¿Cómo se obtiene \fx{Hi} a partir de \fx{hi}?
      \end{enumerate}
    \end{column}
    \begin{column}{0.50\textwidth}
      \begin{enumerate}\setlength{\itemsep}{0.16cm}
        \setcounter{enumi}{10}
        \item ¿Cómo se fuerza el histograma a usar nuestras clases?
        \item ¿Cómo se comparan medias por zona en una línea?
        \item ¿Qué gráfico usa para \fx{zona}? ¿Y para \fx{valor}?
        \item ¿Qué devuelve \fx{boxplot.stats(valor)\$out}?
        \item ¿Por qué \fx{mean(valor)} no coincide con la media agrupada?
        \item ¿Cuál de las dos reportaría y por qué?
        \item ¿Qué significa que \fx{mean()} devuelva \fx{NA}?
        \item ¿Por qué \fx{quantile()} y \fx{fivenum()} difieren?
        \item ¿Por qué una cualitativa debe ser \fx{factor} y no
              \fx{character}?
        \item ¿Cuál es la prueba de que el script está bien hecho?
      \end{enumerate}
    \end{column}
  \end{columns}
\end{frame}

\begin{frame}{Respuestas}{Autoevaluación de los ejercicios propuestos}
  \listatiny
  \begin{columns}[T,onlytextwidth]
    \begin{column}{0.47\textwidth}
      \begin{enumerate}\setlength{\itemsep}{0.16cm}
        \item \fx{mean(valor)}
        \item \fx{median(valor)}
        \item \fx{quantile(valor, 0.75)}
        \item \fx{quantile(valor, 0.40)}
        \item De muestra: divide entre $n-1$
        \item \fx{sd(valor)/mean(valor)*100}
        \item \fx{sum(zona == "Norte")} o \fx{table(zona)}
        \item Mínimo, $Q_1$, mediana, media, $Q_3$ y máximo
        \item \fx{right = FALSE} e \fx{include.lowest = TRUE}
        \item \fx{cumsum(hi)}
      \end{enumerate}
    \end{column}
    \begin{column}{0.50\textwidth}
      \begin{enumerate}\setlength{\itemsep}{0.16cm}
        \setcounter{enumi}{10}
        \item \fx{hist(valor, breaks = limites, right = FALSE)}
        \item \fx{tapply(valor, zona, mean)}
        \item \fx{barplot()} e \fx{hist()}, respectivamente
        \item Los valores atípicos; \fx{numeric(0)} si no hay
        \item Porque la agrupada usa marcas de clase
        \item La de datos crudos, si se tienen; y se declara
        \item Que hay datos faltantes: use \fx{na.rm = TRUE}
        \item Usan convenciones distintas de cuartil
        \item Para fijar los niveles y su orden
        \item Que corre completo desde cero en una sesión nueva
      \end{enumerate}
    \end{column}
  \end{columns}
\end{frame}

\begin{frame}{Referencias}
  \footnotesize
  \begin{itemize}
    \item[{[1]}] Monroy Saldívar, S. (2008). \textit{Estadística descriptiva}.
      Instituto Politécnico Nacional.
      \texttt{https://elibro.net/es/lc/uniminuto/titulos/74722}
    \item[{[2]}] Ordóñez Fernández, F. F. \& González Fernández, J. (2021).
      \textit{Estadística descriptiva paso a paso}. Ediciones Pirámide.
      \texttt{https://elibro.net/es/lc/uniminuto/titulos/215449}
    \item[{[3]}] Llinás Solano, H. (2014). \textit{Introducción a la
      estadística con aplicaciones en ciencias sociales}. Universidad del
      Norte. \texttt{https://elibro.net/es/lc/uniminuto/titulos/70062}
    \item[{[4]}] R Core Team (2025). \textit{An Introduction to R}.
      \texttt{https://cran.r-project.org/doc/manuals/R-intro.pdf}
  \end{itemize}

  \vspace{0.25cm}
  \begin{center}
    {\Large\color{uniAzul}\textbf{¿Preguntas?}}
  \end{center}
\end{frame}

\end{document}
